% ---------------------------------------------------------------------------
% Quem recebeu o dinheiro da cota parlamentar?
% SEMUNI 2026 · UnB · Seminário em Cloud da professora Aleteia · 23/09/2026
%
% 35 min de fala + 10 min de perguntas. 20 slides.
% Compilador: pdfLaTeX (é o padrão do Overleaf — não precisa trocar).
%
% Para gerar o PDF com as notas do apresentador ao lado, descomente a unica
% linha comentada que sobrou logo abaixo do \usepackage{semuni}, e recompile.
% Depois comente de novo: o PDF com notas sai com o dobro da largura e nao
% serve para projetar.
% ---------------------------------------------------------------------------
\documentclass[aspectratio=169,11pt]{beamer}
\usepackage[utf8]{inputenc}
\usepackage[brazil]{babel}
\usepackage{semuni}

% \setbeameroption{show notes on second screen=right}

\title{Quem recebeu o dinheiro da cota parlamentar?}
\author{Ítalo Vinícius}
\date{23 de setembro de 2026}

\begin{document}

% ======================================================== 01 · minuto 0
\begin{frame}
  \rotulo{semuni 2026 · seminário em cloud}
  \titulao{Quem recebeu o dinheiro da cota parlamentar?}
  \vspace{5mm}
  \corpo{Sete anos de nota fiscal, 1,56 milhão de linhas, domínio público. Vou
  responder três vezes nos próximos 35 minutos. As duas primeiras respostas
  estão erradas, e nenhuma delas dá erro.}
  \note{Não se apresente: a professora já leu a minibio. Comece pela pergunta. Diga que você vai responder três vezes e errar duas.}
\end{frame}

% ======================================================== 02 · minuto 1
\begin{frame}
  \rotulo{o dado}
  \subtitulo{Uma nota fiscal por linha, sete arquivos.}
  \vspace{4mm}
  \corpo{A \destaque{CEAP} é a verba mensal de cada deputado. Ele gasta, presta
  contas, e a Câmara publica o CSV inteiro --- um por ano, desde bem antes de
  2019.}
  \vspace{3mm}
  \corpo{\cod{Ano-2019.csv} \ldots \cod{Ano-2025.csv} \quad · \quad 538 MB
  \quad · \quad 32 colunas \quad · \quad \destaque{1.560.019 lançamentos}}
  \note{CEAP = Cota para o Exercício da Atividade Parlamentar. Verba mensal, prestação de contas, publicação obrigatória. Não é vazamento: é transparência ativa.}
\end{frame}

% ======================================================== 03 · minuto 3
\begin{frame}
  \rotulo{o jeito que a gente começa}
  \vspace{2mm}
  \begin{terminal}[analise.py]
  \begin{tela}
  \comentario{import} pandas \comentario{as} pd\\[1.2mm]
  df = pd.read\_csv(\ambar{``Ano-2025.csv''}, sep=\ambar{``;''}, encoding=\ambar{``latin-1''})\\[1.2mm]
  top = (df.groupby(\ambar{``txtFornecedor''})[\ambar{``vlrDocumento''}]\\
  \hphantom{top = (df}.sum()\\
  \hphantom{top = (df}.sort\_values(ascending=\comentario{False})\\
  \hphantom{top = (df}.head(\saida{10}))\\[1.2mm]
  \comentario{print}(top)
  \end{tela}
  \end{terminal}
  \onslide<2->{\legenda{Oito linhas, 40 segundos, nenhuma dependência além do
  pandas. Eu escrevo assim, você escreve assim, e na maior parte das vezes está
  tudo bem.}}
  \note{Mostre o script sem ironia. Esta é a forma correta de começar uma análise: é barata, é rápida, e responde. O problema aparece depois, e não é culpa de quem escreveu.}
\end{frame}

% ======================================================== 04 · minuto 5
\begin{frame}
  \rotulo{a saída, nos sete anos}
  \vspace{2mm}
  \begin{terminal}[python analise.py]
  \begin{tela}
  \vermelho{TAM} \hfill R\$ 116.350.000\\
  GOL \hfill R\$ \phantom{1}52.850.000\\
  \vermelho{Cia Aérea - TAM} \hfill R\$ \phantom{1}43.430.000\\
  AZUL \hfill R\$ \phantom{1}37.530.000\\
  Cia Aérea - GOL \hfill R\$ \phantom{1}35.240.000\\
  PANTANAL VEÍCULOS \hfill R\$ \phantom{1}18.780.000
  \end{tela}
  \end{terminal}
  \onslide<2->{\carimbo{ERRADO}}
  \onslide<2->{\legenda{Primeiro e terceiro lugar são a mesma companhia aérea,
  escrita de dois jeitos. Segundo e quinto também. O script rodou até o fim, sem
  aviso nenhum.}}
  \note{Deixe a tabela na tela em silêncio. Espere alguém na plateia ver sozinho. Só então avance o carimbo.}
\end{frame}

% ======================================================== 05 · minuto 6
\begin{frame}
  \rotulo{dois problemas, empilhados}
  \subtitulo{Nem a soma está certa, nem a pergunta estava bem feita.}
  \vspace{4mm}
  \corpo{\destaque{A soma} quebrou porque o nome do fornecedor é texto livre: a
  mesma empresa entra com grafias diferentes e vira duas linhas.}
  \vspace{3mm}
  \corpo{\destaque{A pergunta} quebrou porque companhia aérea não é fornecedora
  de deputado. É o \destaque{SIGEPA}, o sistema de passagens da própria Câmara:
  o bilhete sai por lá e a companhia cai na planilha.}
  \note{SIGEPA é o sistema de passagens da Câmara. Dois problemas diferentes no mesmo pódio.}
\end{frame}

% ======================================================== 06 · minuto 8
\begin{frame}
  \rotulo{e o dado não fica parado}
  \subtitulo{Quantos lançamentos não têm CNPJ, ano a ano.}
  \vspace{3mm}
  \begin{terminal}[spark-sql --- silver, agrupado por ano]
  \begin{tela}
  ANO \hfill LANÇAMENTOS \hfill SEM CNPJ \hfill \%\\[1mm]
  2019 \hfill 289.830 \hfill 4.448 \hfill \saida{1,5\%}\\
  2020 \hfill 167.132 \hfill 17.234 \hfill 10,3\%\\
  2021 \hfill 218.765 \hfill 43.487 \hfill 19,9\%\\
  2022 \hfill 209.556 \hfill 46.992 \hfill 22,4\%\\
  2023 \hfill 232.745 \hfill 56.943 \hfill 24,5\%\\
  2024 \hfill 232.925 \hfill 58.164 \hfill \vermelho{25,0\%}\\
  2025 \hfill 209.066 \hfill 37.934 \hfill 18,1\%
  \end{tela}
  \end{terminal}
  \onslide<2->{\legenda{A Câmara foi movendo o voo de ``fornecedor com CNPJ'' para
  ``SIGEPA, sem CNPJ''. Quem escreveu o script em 2019 não escreveu nada errado.
  O chão é que andou.}}
  \note{Slide central e o mais difícil de improvisar. Vá devagar. A coluna da direita sobe 16 vezes em cinco anos, e nada no script mudou.}
\end{frame}

% ======================================================== 07 · minuto 10
\begin{frame}
  \frase{Um pipeline não erra menos que um script.\\ Ele erra num lugar onde dá
  para ver.}
  \vspace{6mm}
  \corpo{O que vem agora são quatro erros que eu cometi montando esta aula, e o
  que exatamente pegou cada um deles.}
  \note{Frase-chave. Diga devagar e deixe respirar. Não é um ataque ao script: é a diferença entre errar com rastro e errar sem.}
\end{frame}

% ======================================================== 08 · minuto 12
\begin{frame}
  \rotulo{defeito que não é defeito}
  \subtitulo{56.912 linhas com valor negativo.}
  \vspace{4mm}
  \corpo{São estornos: dinheiro devolvido. Somados junto com o resto, abatem o
  gasto de quem não gastou.}
  \vspace{3mm}
  \onslide<2->{\corpo{Apagar resolveria a soma e perderia a informação. Na silver
  eles viram uma coluna \cod{eh\_estorno}, ficam onde estão, e cada pergunta
  decide se conta com eles. A gold, que responde ``quanto saiu'', conta sem.}}
  \note{Estorno = devolução, valor líquido negativo. A decisão de tirar da soma é legítima; escondê-la não é.}
\end{frame}

% ======================================================== 09 · minuto 14
\begin{frame}
  \rotulo{a mesma empresa, setenta vezes}
  \subtitulo{Setenta fornecedores que são um posto só.}
  \vspace{4mm}
\begin{tikzpicture}
    \foreach \i in {0,...,69}{
      \pgfmathsetmacro{\lin}{int(\i/14)}
      \pgfmathsetmacro{\col}{mod(\i,14)}
      \pgfmathsetmacro{\lar}{Mod(\i,7)==0 ? 7.6 : (Mod(\i,5)==0 ? 3.6 : 5.2)}
      \fill[tinta!40!papel,rounded corners=.3pt]
        (\col*9.2mm,-\lin*3.4mm) rectangle ++(\lar mm,1.5mm);
    }
  \end{tikzpicture}
  \vspace{4mm}
  \onslide<2->{%
    \begin{tikzpicture}[baseline=0pt]
      \fill[verde,rounded corners=.5pt] (0,0) rectangle ++(42mm,2.6mm);
      \node[anchor=west,font=\ttfamily\bfseries\fontsize{8.6}{10}\selectfont,text=tinta]
        at (44mm,1.3mm) {CASCOL --- R\$ 5.522.621 nos sete anos --- uma raiz de CNPJ};
    \end{tikzpicture}%
  }
  \vspace{4mm}
  \onslide<3->{\corpo{Nenhuma grafia isolada era grande o bastante para aparecer
  no pódio. Juntas, entram em décimo.}}
  \note{Cada barrinha é uma grafia diferente do mesmo posto. A raiz do CNPJ — os 8 primeiros dígitos — é o que une, porque filial muda os 4 do meio.}
\end{frame}

% ======================================================== 10 · minuto 16
\begin{frame}
  \rotulo{uma assinatura de ChatGPT, dois universos}
  \vspace{2mm}
  \begin{columns}[t,onlytextwidth]
    \begin{column}{.48\textwidth}
      \begin{lado}{2023}
        {\ttfamily\bfseries\fontsize{11}{13}\selectfont\color{tinta}00000000000010\par}\vspace{1.6mm}
        {\fontsize{8}{10.4}\selectfont\color{\tintaSessenta}\destaque{ChatGPT Plus Subscription}\\
        R\$ 98,93, com R\$ 56,06 de glosa\\ categoria: manutenção de escritório\par}
      \end{lado}
    \end{column}
    \begin{column}{.48\textwidth}
      \begin{lado}{2025}
        {\ttfamily\bfseries\fontsize{11}{13}\selectfont\color{laranja}62531071000178\par}\vspace{1.6mm}
        {\fontsize{8}{10.4}\selectfont\color{\tintaSessenta}\destaque{chatGPT}\\
        R\$ 999,90\\ categoria: manutenção de escritório\par}
      \end{lado}
    \end{column}
  \end{columns}
  \vspace{5mm}
  \onslide<2->{\fecho{O documento da esquerda tem 14 dígitos e passa em qualquer
  validação de CNPJ. Ele é uma gaveta interna da Câmara, e naquele mesmo ano a
  mesma assinatura também foi lançada como \destaque{assinatura de publicações},
  por R\$ 105,03.}}
  \note{O caso que junta todos os defeitos numa linha só, e é recente o bastante para a plateia rir.}
\end{frame}

% ======================================================== 11 · minuto 18
\begin{frame}
  \rotulo{onde o conserto vai morar}
  \vspace{1mm}
  \begin{columns}[t,onlytextwidth]
    \begin{column}{.315\textwidth}
      \begin{lado}{camada 1}
        {\bfseries\fontsize{12}{14}\selectfont\color{tinta}bronze\par}\vspace{1.4mm}
        {\fontsize{8}{10.4}\selectfont\color{\tintaSessenta}Como a Câmara
        publicou. Tudo texto, nada consertado. 1.560.019 linhas.\par}
      \end{lado}
    \end{column}
    \begin{column}{.315\textwidth}
      \begin{lado}{camada 2}
        {\bfseries\fontsize{12}{14}\selectfont\color{tinta}silver\par}\vspace{1.4mm}
        {\fontsize{8}{10.4}\selectfont\color{\tintaSessenta}Tipos certos e sete
        colunas de defeito marcado. Nenhuma linha some.\par}
      \end{lado}
    \end{column}
    \begin{column}{.315\textwidth}
      \begin{lado}{camada 3}
        {\bfseries\fontsize{12}{14}\selectfont\color{tinta}gold\par}\vspace{1.4mm}
        {\fontsize{8}{10.4}\selectfont\color{\tintaSessenta}Uma tabela por
        pergunta. 55.672 empresas, somadas pela raiz do CNPJ.\par}
      \end{lado}
    \end{column}
  \end{columns}
  \vspace{5mm}
  \onslide<2->{\fecho{Guardar o bronze cru custa 5 GB e parece desperdício até a
  primeira vez que a regra de limpeza está errada. Aí ele é o que permite refazer
  sem pedir o arquivo de novo --- e a Câmara \destaque{republica os CSVs com
  correções}, então o arquivo de hoje não é o de ontem.}}
  \note{Agora a estrutura, e só agora. O bronze existe porque a sua regra vai estar errada uma hora.}
\end{frame}

% ======================================================== 12 · minuto 20
\begin{frame}
  \rotulo{a mesma pergunta, nas duas pontas}
  \vspace{2mm}
  \begin{terminal}[spark-sql --- container spark-master]
  \begin{tela}
  \sql{SELECT txtFornecedor, sum(vlrDocumento) t}\\
  \hphantom{spark-sql> }FROM delta.`s3a://lake/\vermelho{bronze}` WHERE numAno=2025 GROUP BY 1 ORDER BY 2 DESC LIMIT 3;\\
  \vermelho{TAM 16.402.463,26 \quad GOL 6.052.550,88 \quad AZUL 4.696.245,04}\\[1.6mm]
  \sql{SELECT fornecedor, total}\\
  \hphantom{spark-sql> }FROM fornecedores WHERE ano\_ref=2025 ORDER BY total DESC LIMIT 3;\\
  \saida{FACEBOOK 3.251.951,87 \quad PANTANAL 3.227.314,77 \quad VIVO 1.591.715,84}
  \end{tela}
  \end{terminal}
  \onslide<2->{\legenda{Mesmo arquivo, mesma pergunta, nenhum nome repetido entre
  as duas listas.}}
  \note{A mesma pergunta nas duas pontas. Aponte que a de cima é o script do slide 3, só que em SQL.}
\end{frame}

% ======================================================== 13 · minuto 22
\begin{frame}
  \rotulo{a resposta, em 2025}
  \vspace{4mm}
  \gigante{verde}{R\$ 3,25 mi}
  \vspace{1.5mm}
  \legendanum{FACEBOOK SERVIÇOS ONLINE DO BRASIL}
  \vspace{6mm}
  \onslide<2->{\corpo{Nos sete anos somados a lista volta a começar com companhia
  aérea, e dessa vez está certo: até 2022 a Câmara lançava o voo \destaque{com}
  CNPJ. A resposta depende do ano porque o dado mudou de forma, e é por isso que
  o ano é uma partição e não um filtro solto no meio de um script.}}
  \note{Clímax. Silêncio antes do fragmento. E seja honesto: nos sete anos a lista muda, e o motivo é o slide 6.}
\end{frame}

% ======================================================== 14 · minuto 24
\begin{frame}
  \rotulo{quatro erros meus, e o que pegou cada um}
  \vspace{2mm}
  \begin{terminal}[o que aconteceu montando esta aula]
  \begin{tela}
  O ERRO \hfill O QUE PEGOU\\[1.4mm]
  regra de glosa acusou 53.112 linhas boas \hfill \saida{o relatório conta, não filtra}\\
  a mesma empresa com nome diferente a cada execução \hfill \saida{rodar duas vezes e comparar}\\
  soma saindo como 4.99577E7 na tela \hfill \saida{tipo decimal na gold}\\
  \vermelho{R\$ 15,1 mi de um fornecedor que não existe} \hfill \saida{consultar a gold antes de crer}
  \end{tela}
  \end{terminal}
  \onslide<2->{\legenda{O quarto ia para este slide. Ramal, celular funcional e
  Correios compartilham a raiz \cod{00000000}, viraram uma empresa só e
  apareceram em quarto lugar. Os catorze dígitos batem, inclusive os
  verificadores. Formato válido e significado válido são coisas diferentes.}}
  \note{O coração da aula. São quatro erros meus, não hipotéticos. Se estiver atrasado, corte o quarto.}
\end{frame}

% ======================================================== 15 · minuto 27
\begin{frame}
  \rotulo{o erro continua consultável}
  \vspace{2mm}
  \begin{terminal}[spark-sql --- DESCRIBE HISTORY]
  \begin{tela}
  VERSÃO \hfill OPERAÇÃO \hfill LINHAS \hfill O QUE ERA\\[1.4mm]
  0 \hfill WRITE \hfill \vermelho{209.079} \hfill sem multiLine --- 13 linhas partidas\\
  1 \hfill WRITE \hfill \saida{209.066} \hfill com multiLine --- correto\\[2.4mm]
  \sql{SELECT count(*) FROM delta.`s3a://lake/bronze` VERSION AS OF 0;}\\
  \vermelho{209079}
  \end{tela}
  \end{terminal}
  \onslide<2->{\legenda{Treze linhas têm quebra de linha dentro de um campo entre
  aspas, e dois motores honestos discordam sobre isso por padrão. O que a tabela
  Delta acrescenta é que a contagem errada não foi sobrescrita: ela virou a
  versão 0, e dá para voltar nela agora.}}
  \note{Aconteceu de verdade: DuckDB contou 209.066 e Spark contou 209.079. A v0 guardou o erro. Consulte ao vivo.}
\end{frame}

% ======================================================== 16 · minuto 29
\begin{frame}
  \rotulo{o que impede a volta}
  \subtitulo{21 testes, e dois deles têm nome de cicatriz.}
  \vspace{3mm}
  \begin{terminal}[pytest /opt/testes]
  \begin{tela}
  \saida{PASSED} test\_estorno\_nao\_conta\_como\_glosa\_maior\\
  \saida{PASSED} test\_raiz\_zerada\_nao\_e\_empresa\\
  \saida{PASSED} test\_raiz\_de\_verdade\_nao\_e\_confundida\\
  \comentario{...}\\[1.2mm]
  \saida{21 passed} in 11.09s
  \end{tela}
  \end{terminal}
  \onslide<2->{\legenda{Rodam sem MinIO e sem cluster, numa SparkSession local,
  em 11 segundos. O que eles testam não é o Spark: é se a regra de negócio ainda
  significa o que eu quis dizer.}}
  \note{Dois destes testes existem porque os bugs aconteceram. Teste em engenharia de dados é sobre a regra, não sobre o framework.}
\end{frame}

% ======================================================== 17 · minuto 30
\begin{frame}
  \rotulo{a sala de máquinas --- tudo em containers}
  % o frame e [t], entao o diagrama encostaria no rotulo e deixaria um vazio de
  % ~20mm sobre o rodape. Os dois \vfill centram o desenho no espaco que sobra.
  \vfill
  \begin{center}
  \begin{tikzpicture}[x=1mm,y=1mm]
    % processamento
    \draw[slab] (0mm,3mm) rectangle (142mm,-13mm);
    \node[ts] at (2.5mm,1mm) {PROCESSAMENTO · CONTAINERS spark};
    \bloco{2.5}{-2}{30}{cx}{spark-master}{REPARTE O TRABALHO}
    \bloco{35}{-2}{25}{cx}{worker 1}{2 CORES · 2 GB}
    \bloco{62.5}{-2}{25}{cx}{worker 2}{2 CORES · 2 GB}
    \blocoq{90}{-2}{25}{worker 3}{SOBE AO VIVO}

    % o cano
    \draw[cano] (60mm,-15mm) -- (60mm,-22mm);
    \draw[cano] (54mm,-22mm) -- (54mm,-15mm);
    \node[ts] at (66mm,-15.5mm) {s3a:// --- A MESMA API DO AMAZON S3};
    \node[tp] at (66mm,-19.2mm) {TROCA O ENDPOINT E O MESMO CÓDIGO RODA NA AWS};

    % armazenamento
    \draw[slab] (0mm,-24mm) rectangle (142mm,-40mm);
    \node[ts] at (2.5mm,-26mm) {ARMAZENAMENTO · CONTAINER minio};
    \bloco{2.5}{-29}{30}{cx}{MinIO}{BUCKET lake}
    \bloco{35}{-29}{25}{cx}{bronze}{1.560.019 LINHAS}
    \bloco{62.5}{-29}{25}{cx}{silver}{DEFEITOS MARCADOS}
    \bloco{90}{-29}{25}{cx}{gold}{55.672 EMPRESAS}

    \node[tn] at (2.5mm,-44mm)
      {o Spark não guarda o dado. derrube os três containers dele e o lake continua inteiro.};
  \end{tikzpicture}
  \end{center}
  \vfill
  \note{O worker 3 nasce tracejado porque vai subir de verdade: docker compose up -d --scale spark-worker=3, com o Spark UI aberto noutra aba.}
\end{frame}

% ======================================================== 18 · minuto 31
\begin{frame}
[t]
  \rotulo{o mesmo desenho, com o nome que ele tem na nuvem}
  \vspace{2mm}
  \begin{center}
  \begin{tikzpicture}[x=1mm,y=1mm]
    \draw[slab] (0mm,3mm) rectangle (142mm,-12mm);
    \node[ts] at (2.5mm,1mm) {PROCESSAMENTO};
    \bloco{2.5}{-2}{32}{cx}{Apache Spark}{NO SEU DOCKER}
    \draw[seta] (36.5mm,-6.5mm) -- (41.5mm,-6.5mm);
    \bloco{43.5}{-2}{96}{cx}{EMR · Dataproc · Databricks · Synapse}{O MESMO APACHE SPARK, OPERADO POR OUTRA PESSOA}

    \draw[slab] (0mm,-17mm) rectangle (142mm,-32mm);
    \node[ts] at (2.5mm,-19mm) {ARMAZENAMENTO};
    \bloco{2.5}{-22}{32}{cx}{MinIO}{NO SEU DOCKER}
    \draw[seta] (36.5mm,-26.5mm) -- (41.5mm,-26.5mm);
    \bloco{43.5}{-22}{96}{cx}{Amazon S3 · Google Cloud Storage · Azure Blob}{A MESMA API. VOCÊ TROCA UMA LINHA DE CONFIGURAÇÃO.}

    \node[tn] at (2.5mm,-37mm)
      {Delta Lake é o mesmo formato dos dois lados. Ele é aberto, e não muda quando você paga.};
  \end{tikzpicture}
  \end{center}
  \note{Tom neutro. É o mesmo software rodando em dois lugares, e o código não muda. Não deprecie nenhum lado.}
\end{frame}

% ======================================================== 19 · minuto 32
\begin{frame}
  \rotulo{quando cada um faz sentido}
  \vspace{1mm}
  \begin{columns}[t,onlytextwidth]
    \begin{column}{.48\textwidth}
      \begin{lado}{no seu docker}
        {\bfseries\fontsize{11}{13}\selectfont\color{tinta}Bom quando\dots\par}
        \vspace{1mm}
        \begin{beneficios}
          \item você está aprendendo e quer quebrar sem medo
          \item o dado cabe confortavelmente numa máquina
          \item o ciclo curto importa: subiu, errou, refez em minutos
          \item o dado não pode sair da sua rede
        \end{beneficios}
      \end{lado}
    \end{column}
    \begin{column}{.48\textwidth}
      \begin{lado}{na nuvem}
        {\bfseries\fontsize{11}{13}\selectfont\color{tinta}Bom quando\dots\par}
        \vspace{1mm}
        \begin{beneficios}
          \item o dado não cabe mais numa máquina só
          \item você precisa de 200 máquinas por 20 minutos
          \item o time é pequeno e o plantão não pode ser você
          \item outras equipes precisam do mesmo dado, com governança
        \end{beneficios}
      \end{lado}
    \end{column}
  \end{columns}
  \onslide<2->{\fecho{É o \destaque{mesmo Apache Spark} e o \destaque{mesmo Delta
  Lake} dos dois lados. O que muda é quem opera a infraestrutura.}}
  \note{Os dois lados têm bolinha verde. A escolha é de contexto, não de virtude.}
\end{frame}

% ======================================================== 20 · minuto 33
\begin{frame}
  \rotulo{para levar para casa}
  \subtitulo{O script do slide 3 continua certo.\\ Ele só precisava de um lugar
  para morar.}
  \vspace{5mm}
  \corpo{Tudo que você viu é software aberto --- MinIO, Apache Spark, Delta Lake,
  Docker --- e um CSV público da Câmara. Roda no seu laptop hoje e sobe na nuvem
  no dia em que o dado crescer, sem reescrever a regra.}
  \vspace{6mm}
  {\ttfamily\fontsize{10}{12}\selectfont\color{tinta}%
   make subir \quad · \quad make pipeline\par}
  \note{Fecho curto. O repositório sobe em dois comandos. Depois abra para as perguntas — são 10 minutos.}
\end{frame}

\end{document}
